\chapter{Gravitasi Semesta dan Berat}\label{ch:g10:universal-gravitation}

Jatuhkanlah sebuah apel: apel itu jatuh. Bulan, enam puluh jari-jari
Bumi di atas sana, tampaknya tidak pernah jatuh. Wawasan besar Newton
adalah bahwa Bulan \emph{sedang} jatuh --- kira-kira 1.4 milimeter tiap
sekon, tepat sebanyak yang diramalkan tarikan yang melemah menurut
kuadrat jaraknya. Satu rumus, dengan satu tetapan semesta, mengatur
apel, Bulan, pasang surut dan planet. Jilid pertama melukiskan \omterm{def:g10:universal-gravitation:force}{gravitasi}
dengan kata-kata; bab ini mengukurnya.

\section{Hukum gravitasi semesta}

\begin{definition}[Gaya gravitasi]\label{def:g10:universal-gravitation:force}
Dua benda apa pun saling menarik, dari bahan apa pun keduanya terbuat
dan seberapa jauh pun keduanya terpisah. Tarikan itulah \emph{gaya
gravitasi}\index{gaya gravitasi}, dan gejalanya disebut
\emph{gravitasi}\index{gravitasi}. Gaya itu tidak memerlukan sentuhan,
tidak memerlukan tali dan tidak memerlukan medium, dan gaya itu bersifat
\emph{timbal balik}: masing-masing benda menarik benda yang lain.
\end{definition}

\begin{theorem}[Hukum gravitasi semesta]\label{thm:g10:universal-gravitation:law}
Dua benda bermassa $m_1$ dan $m_2$ (dalam \unit{kg}), yang pusatnya
terpisah sejauh $d$ (dalam \unit{m}), saling menarik dengan dua gaya
yang terarah sepanjang garis penghubung kedua pusatnya, masing-masing
menuju benda yang lain, dan besarnya sama, yaitu
\[
F = G\,\frac{m_1 m_2}{d^2}
\qquad \text{(dalam \unit{N})},
\]
dengan $G$ adalah tetapan yang sama bagi setiap pasang benda di alam
semesta.
\end{theorem}
\admitted

\begin{remark}[Dari mana hukum ini datang]
Newton menyaring hukum ini dari gerak planet yang teramati (hukum
Kepler); penurunannya yang jujur dikerjakan pada jilid Tahun ke-1.
\Cref{ch:g12:satellites-kepler} sudah akan memakai hukum ini untuk
\omterm{def:g10:universal-gravitation:satellite}{satelit} dan planet.
\end{remark}

\begin{definition}[Tetapan gravitasi]\label{def:g10:universal-gravitation:constant}
Tetapan
\[
G = \qty{6.67e-11}{N.m^2/kg^2}
\]
disebut \emph{tetapan gravitasi}\index{tetapan gravitasi}. Nilainya yang
teramat kecil adalah alasan mengapa \omterm{def:g10:universal-gravitation:force}{gravitasi} antara benda sehari-hari
tidak terasa: hanya apabila sekurang-kurangnya satu massanya berukuran
astronomi barulah $G m_1 m_2 / d^2$ berarti sesuatu.
\end{definition}

\begin{omfigure}
\begin{tikzpicture}[scale=1.0]
  \draw[omExo, dashed] (0,0) -- (5.2,0);
  \fill[omDef!25] (0,0) circle (0.45);
  \draw[omDef, thick] (0,0) circle (0.45);
  \fill[omDef!25] (5.2,0) circle (0.28);
  \draw[omDef, thick] (5.2,0) circle (0.28);
  \node[font=\small] at (0,0) {$m_1$};
  \node[font=\small] at (5.2,0) {$m_2$};
  \draw[-Stealth, omThm, very thick] (0.45,0) -- (1.55,0)
    node[midway, above, font=\small] {$\vect F_{2/1}$};
  \draw[-Stealth, omThm, very thick] (4.92,0) -- (3.82,0)
    node[midway, above, font=\small] {$\vect F_{1/2}$};
  \draw[Stealth-Stealth, thin] (0,-0.75) -- (5.2,-0.75)
    node[midway, below, font=\small] {$d$};
\end{tikzpicture}

{\small Dua gaya sebuah pasangan \omterm{def:g10:universal-gravitation:force}{gravitasi}: sepanjang garis
penghubung pusatnya, menuju benda yang lain, dengan besar yang
\emph{sama} sekalipun massanya jauh berbeda,
$F = G m_1 m_2 / d^2$.}
\end{omfigure}

\begin{remark}[Membaca rumusnya]\label{rem:g10:universal-gravitation:reading}
Ada empat ciri yang patut diperhatikan.
\begin{itemize}
  \item \emph{Timbal balik dan sama besar.} Bumi menarikmu dengan
        beratmu --- dan kamu menarik Bumi dengan gaya yang persis sama.
        (Bumi hanya menanggapinya lebih sedikit: massanya sekitar
        $\num{e23}$ kali lebih besar.)
  \item \emph{Kebalikan kuadrat.} Menggandakan $d$ membagi gayanya
        dengan $4$; menjauh sepuluh kali lipat membaginya dengan $100$.
  \item \emph{Jarak antara pusat.} Untuk benda berbentuk bola ---
        bintang, planet, peluru meriam --- $d$ diukur dari pusat ke
        pusat, bukan dari permukaan ke permukaan.
  \item \emph{Kedua massa berperan.} Menggandakan salah satu massanya
        menggandakan gayanya; rumusnya setangkup dalam $m_1$ dan $m_2$.
\end{itemize}
\end{remark}

\begin{example}[Dua sahabat]\label{ex:g10:universal-gravitation:friends}
Dua sahabat bermassa \qty{70}{kg} berdiri dengan pusatnya terpisah
\qty{1.0}{m}:
\[
F = \num{6.67e-11} \times \frac{70 \times 70}{1.0^2}
\approx \qty{3.3e-7}{N}.
\]
\omterm{def:g10:universal-gravitation:weight}{Berat} masing-masing sahabat itu sekitar \qty{6.9e2}{N} --- dua
\emph{miliar} kali lebih besar. Tarikan timbal balik mereka nyata,
tetapi mustahil dirasakan.
\end{example}

\begin{example}[Bumi menahan Bulan]\label{ex:g10:universal-gravitation:earthmoon}
Dengan $m_1 = \qty{5.97e24}{kg}$ (Bumi), $m_2 = \qty{7.35e22}{kg}$
(Bulan) dan $d = \qty{3.84e8}{m}$:
\[
F = \num{6.67e-11} \times
\frac{\num{5.97e24} \times \num{7.35e22}}{(\num{3.84e8})^2}
\approx \qty{2.0e20}{N}.
\]
Dan Bulan menarik Bumi kembali dengan \qty{2.0e20}{N} yang sama ---
tarikan yang dijawab lautan dua kali sehari, sebagai pasang surut.
\end{example}

\begin{method}[Menghitung gaya gravitasi]\label{met:g10:universal-gravitation:compute}
\begin{enumerate}
  \item Ubahlah massanya ke \omterm{def:g10:orders-of-magnitude:unit}{kilogram} dan jarak pusat ke pusatnya ke
        \omterm{def:g10:orders-of-magnitude:unit}{meter}.
  \item Terapkan $F = G m_1 m_2 / d^2$.
  \item Periksalah pangkat sepuluhnya terpisah dari angkanya, seperti
        pada \cref{ch:g10:orders-of-magnitude}: \omterm{def:g10:universal-gravitation:force}{gaya gravitasi}
        berkisar dari $\num{e-7}$ newton (dua orang yang terpisah satu
        \omterm{def:g10:orders-of-magnitude:unit}{meter}) sampai $\num{e22}$ newton (Matahari pada Bumi),
        sehingga \omterm{def:g10:orders-of-magnitude:scinot}{eksponen} yang salah tempat mudah tertangkap.
\end{enumerate}
\end{method}

\section{Berat dan kuat medan gravitasi}

\begin{definition}[Berat]\label{def:g10:universal-gravitation:weight}
\emph{Berat}\index{berat} $\vect P$ sebuah benda di permukaan sebuah
benda langit (planet, bulan, \dots) adalah \omterm{def:g10:universal-gravitation:force}{gaya gravitasi} yang
dikerjakan benda langit itu padanya. Berat menunjuk ke pusat benda
langit tersebut --- arah yang kita sebut
\emph{vertikal}\index{vertikal} --- dan diukur dalam newton.
\end{definition}

\begin{proposition}[Berat di permukaan sebuah benda langit]\label{prop:g10:universal-gravitation:surfaceg}
Di permukaan benda langit berbentuk bola yang bermassa $M$ dan
berjari-jari $R$, sebuah benda bermassa $m$ memiliki \omterm{def:g10:universal-gravitation:weight}{berat}
\[
P = m\,g,
\qquad\text{dengan}\qquad
g = \frac{G M}{R^2}
\]
yang hanya bergantung pada benda langitnya, bukan pada bendanya.
\end{proposition}

\begin{proof}
Pusat bendanya dan pusat benda langitnya terpisah sejauh $R$, sehingga
hukum \omterm{def:g10:universal-gravitation:force}{gravitasi} semesta (\cref{thm:g10:universal-gravitation:law})
memberikan $P = G M m / R^2 = m \times \left(G M / R^2\right)$. Kurungan
itu sama bagi setiap benda di permukaannya: sebutlah ia $g$.
\end{proof}

\begin{definition}[Kuat medan gravitasi]\label{def:g10:universal-gravitation:fieldstrength}
Besaran $g = G M / R^2$, dalam newton per \omterm{def:g10:orders-of-magnitude:unit}{kilogram}, disebut \emph{kuat
medan gravitasi}\index{kuat medan gravitasi} di permukaan benda langit
itu: benda langit tersebut menarik tiap \omterm{def:g10:orders-of-magnitude:unit}{kilogram} dengan $g$ newton. Di
Bumi, $g = \qty{9.81}{N/kg}$.
\end{definition}

\begin{example}[Menghitung $g$]\label{ex:g10:universal-gravitation:gvalues}
Untuk Bumi ($M = \qty{5.97e24}{kg}$, $R = \qty{6.37e6}{m}$):
\[
g = \frac{\num{6.67e-11} \times \num{5.97e24}}{(\num{6.37e6})^2}
= \frac{\num{3.98e14}}{\num{4.06e13}}
\approx \qty{9.81}{N/kg}.
\]
Untuk Bulan ($M = \qty{7.35e22}{kg}$, $R = \qty{1.74e6}{m}$)
perhitungan yang sama memberikan
$g_{\text{Bulan}} \approx \qty{1.62}{N/kg}$ --- enam kali lebih kecil:
massa Bulan jauh lebih ringan, dan ukurannya yang lebih kecil tidak
menutupi kekurangan itu. Untuk Mars,
$g_{\text{Mars}} \approx \qty{3.73}{N/kg}$.
\end{example}

\begin{omfigure}
\begin{tikzpicture}
\begin{axis}[omaxis, axis x line=bottom, axis y line=left,
  width=8.8cm, height=5.4cm, ybar, bar width=15pt,
  symbolic x coords={Bulan, Mars, Bumi, Jupiter}, xtick=data,
  ylabel={$g$ (\unit{N/kg})}, ymin=0, ymax=29,
  nodes near coords,
  every node near coord/.append style={font=\small},
  enlarge x limits=0.16]
\addplot[draw=omDef, fill=omDef!25] coordinates
  {(Bulan,1.62) (Mars,3.73) (Bumi,9.81) (Jupiter,24.8)};
\end{axis}
\end{tikzpicture}

{\small \omterm{def:g10:universal-gravitation:fieldstrength}{Kuat medan gravitasi} di permukaan empat dunia: sekantong gula
\qty{1}{kg} yang sama \omterm{def:g10:universal-gravitation:weight}{beratnya} \qty{1.6}{N} di Bulan dan \qty{25}{N} di
Jupiter.}
\end{omfigure}

\begin{omfigure}
\begin{tikzpicture}[scale=1.0]
  \fill[omExo!15] (0,0) circle (1.5);
  \draw[omExo, thick] (0,0) circle (1.5);
  \fill (0,0) circle (1.5pt) node[below, font=\small] {$O$};
  \fill (90:1.5) circle (1.8pt) node[above, font=\small] {$A$};
  \draw[-Stealth, omThm, thick] (90:1.5) -- (90:0.8)
    node[pos=0.6, right, font=\small] {$\vect P_A$};
  \fill (200:1.5) circle (1.8pt) node[left, font=\small] {$B$};
  \draw[-Stealth, omThm, thick] (200:1.5) -- (200:0.8)
    node[pos=0.6, below, font=\small] {$\vect P_B$};
  \fill (330:1.5) circle (1.8pt) node[right, font=\small] {$C$};
  \draw[-Stealth, omThm, thick] (330:1.5) -- (330:0.8)
    node[pos=0.6, above left, font=\small] {$\vect P_C$};
\end{tikzpicture}

{\small ``Bawah'' adalah gagasan yang setempat: di setiap titik bola
Bumi vektor \omterm{def:g10:universal-gravitation:weight}{beratnya} menunjuk ke pusat $O$. Dua pejalan di dua kutub
yang berhadapan sama-sama berdiri tegak --- kaki mereka saling
berhadapan.}
\end{omfigure}

\begin{remark}[Mengapa antariksawan memantul-mantul]\label{rem:g10:universal-gravitation:bounce}
Seorang antariksawan bermassa \qty{120}{kg} (tubuh beserta pakaiannya)
\omterm{def:g10:universal-gravitation:weight}{beratnya} $120 \times 9.81 \approx \qty{1180}{N}$ di Bumi tetapi hanya
$120 \times 1.62 \approx \qty{194}{N}$ di Bulan. Otot yang terlatih di
bawah \qty{1180}{N} tiba-tiba memikul seperenam bebannya: dari sinilah
lahir langkah melompat yang lambat dan tersohor di Bulan itu. Massanya
--- dan bersamanya usaha yang diperlukan untuk \emph{berhenti} atau
\emph{berbelok} --- tidak berubah, dan itulah sebabnya antariksawan di
Bulan begitu sering terjatuh.
\end{remark}

\section{Massa atau berat? Neraca dan neraca pegas}

Bahasa sehari-hari mengatakan sekarung tepung ``\omterm{def:g10:universal-gravitation:weight}{beratnya} tiga kilo''.
Fisika membelah kalimat itu menjadi dua.

\begin{remark}[Massa bukan berat]\label{rem:g10:universal-gravitation:massweight}
\begin{itemize}
  \item \emph{Massa} $m$ mengukur banyaknya materi, dalam \omterm{def:g10:orders-of-magnitude:unit}{kilogram}.
        Massa adalah sifat bendanya sendiri: sama di Bumi, di Bulan,
        atau ketika melayang di ruang angkasa yang jauh.
  \item \emph{\omterm{def:g10:universal-gravitation:weight}{Berat}} $P = mg$ adalah sebuah gaya, dalam newton, yang
        dikerjakan \emph{oleh sebuah benda langit} pada bendanya.
        \omterm{def:g10:universal-gravitation:weight}{Berat} berubah dari benda langit ke benda langit mengikuti $g$
        --- dan memudar jauh dari setiap benda langit.
\end{itemize}
Mencampuradukkan keduanya tidak berbahaya di pasar dan berakibat fatal
di dalam soal fisika: \omterm{def:g10:orders-of-magnitude:unit}{kilogram} bukan \omterm{def:g10:orders-of-magnitude:unit}{satuan} gaya.
\end{remark}

\begin{method}[Neraca dan neraca pegas]\label{met:g10:universal-gravitation:instruments}
\begin{itemize}
  \item \emph{Neraca lengan} membandingkan bendanya dengan massa
        baku: \omterm{def:g10:universal-gravitation:weight}{berat} kedua piringnya berskala dengan $g$ yang sama,
        sehingga perbandingannya mengukur \emph{massa} --- putusan yang
        sama di benda langit mana pun.
  \item \emph{Neraca pegas} mengukur gaya yang meregangkan pegasnya:
        alat itu membaca \emph{\omterm{def:g10:universal-gravitation:weight}{berat}}, dalam newton, dan bacaannya
        berubah dari benda langit ke benda langit. Untuk memperoleh
        kembali massanya, bagilah dengan kuat medan setempat:
        $m = P/g$.
\end{itemize}
\end{method}

\begin{example}[Tepung di Bulan]\label{ex:g10:universal-gravitation:flour}
Sekarung tepung \qty{3.0}{kg} setimbang dengan massa baku \qty{3.0}{kg}
di Bumi \emph{dan} di Bulan. Sebaliknya, neraca pegas membaca
$3.0 \times 9.81 \approx \qty{29}{N}$ di Bumi dan
$3.0 \times 1.62 \approx \qty{4.9}{N}$ di Bulan. Pedagang yang menjual
``per newton'' di Bulan akan menyerahkan tepung enam kali lipat ---
perdagangan, seperti ilmu, berjalan di atas massa.
\end{example}

\section{Jatuh mengitari Bumi: meriam Newton}

Lemparkanlah sebuah batu mendatar: \omterm{def:g10:universal-gravitation:force}{gravitasi} melengkungkan lintasannya
menjadi busur dan batu itu mendarat beberapa \omterm{def:g10:orders-of-magnitude:unit}{meter} jauhnya.
Lemparkanlah lebih cepat: batu itu mendarat lebih jauh. Percobaan
pikiran Newton mendorong gagasan itu sampai batasnya. Bumi berbentuk
bulat, dan permukaannya turun sekitar \qty{5}{m} di bawah garis mendatar
untuk setiap \qty{8}{km} yang ditempuh; benda yang jatuh bebas turun
\qty{4.9}{m} pada sekon pertamanya
(\cref{ch:g12:projectile-motion} akan menetapkannya). Jadi peluru yang
menyerempet Bumi pada \qty{8}{km/s} jatuh persis secepat tanah
melengkung menjauh: peluru itu jatuh selamanya tanpa pernah mendarat.

\begin{definition}[Satelit dan orbit]\label{def:g10:universal-gravitation:satellite}
\emph{Satelit}\index{satelit} sebuah benda langit adalah benda yang
jatuh mengitari benda langit itu tanpa pernah mencapai permukaannya;
lintasan yang diulanginya disebut \emph{orbit}\index{orbit}.
\end{definition}

\begin{omfigure}
\begin{tikzpicture}[scale=1.0]
  \fill[omExo!15] (0,0) circle (1.5);
  \draw[omExo, thick] (0,0) circle (1.5);
  \fill[omExo!60] (-0.22,1.44) -- (0,1.85) -- (0.22,1.44) -- cycle;
  \coordinate (L) at (0,1.85);
  \draw[omThm, thick] (L) .. controls (0.6,1.85) .. (55:1.5);
  \draw[omProp, thick] (L) .. controls (1.1,1.85) and (1.7,1.1) .. (10:1.5);
  \draw[omMeth, thick] (L) .. controls (1.8,1.8) and (2.3,-0.4) .. (-70:1.5);
  \draw[-Stealth, omDef, thick] (0,1.85) arc (90:-250:1.85);
  \node[font=\small, omDef, anchor=east] at (-1.5,1.45) {orbit};
\end{tikzpicture}

{\small Meriam Newton: ditembakkan makin lama makin cepat dari puncak
gunung, bolanya mendarat makin lama makin jauh --- sampai, pada sekitar
\qty{7.9}{km/s}, tanah melengkung menjauh secepat bolanya jatuh, dan
bola itu pun mengorbit.}
\end{omfigure}

\begin{remark}[Bulan adalah batu yang sedang jatuh]\label{rem:g10:universal-gravitation:fallingmoon}
Bulan adalah \omterm{def:g10:universal-gravitation:satellite}{satelit} alami Bumi: ia jatuh ke arah kita sekitar
\qty{1.4}{mm} tiap sekon, dan gerak menyampingnya membawanya berkeliling
sebelum ia sempat mendekat
(\cref{exo:g10:universal-gravitation:15} mengajakmu mengulang
pemeriksaan Newton sendiri atas angka ini). Antariksawan di stasiun yang
mengorbit berada dalam jatuh bebas yang sama dengan wahananya --- itulah,
dan bukan ketiadaan \omterm{def:g10:universal-gravitation:force}{gravitasi}, sebabnya mereka melayang. \omterm{def:g10:universal-gravitation:force}{Gravitasi}
membentuk setiap lintasan, dari busur batu sampai elips planet; kisah
\omterm{def:g10:universal-gravitation:satellite}{orbit} yang terhitung dituturkan pada \cref{ch:g12:satellites-kepler}.
\end{remark}

\begin{notation}[Kartu data]
Kecuali jika sebuah latihan menyatakan lain, pakailah
$G = \qty{6.67e-11}{N.m^2/kg^2}$ dan:
\begin{center}
\small
\begin{tabular}{lccc}
 & massa (\unit{kg}) & jari-jari (\unit{m}) & $g$ (\unit{N/kg}) \\
Bumi & \num{5.97e24} & \num{6.37e6} & 9.81 \\
Bulan & \num{7.35e22} & \num{1.74e6} & 1.62 \\
Mars & \num{6.42e23} & \num{3.39e6} & 3.73 \\
Matahari & \num{1.99e30} & --- & --- \\
\end{tabular}
\end{center}
Jarak Bumi--Bulan: \qty{3.84e8}{m}; jarak Bumi--Matahari:
\qty{1.496e11}{m} (pusat ke pusat).
\end{notation}

\section{Latihan}

\begin{exercise}[$\star$]\label{exo:g10:universal-gravitation:1}
Dua penari bermassa \qty{60}{kg} saling berpegangan dengan pusatnya
terpisah \qty{0.80}{m}.
\begin{enumerate}
  \item Hitunglah tarikan \omterm{def:g10:universal-gravitation:force}{gravitasi} timbal balik keduanya.
  \item Bandingkan dengan \omterm{def:g10:universal-gravitation:weight}{berat} seekor nyamuk \qty{2.5}{mg}.
\end{enumerate}
\end{exercise}

\begin{exercise}[$\star$]\label{exo:g10:universal-gravitation:2}
Seorang pelajar bermassa \qty{55}{kg}.
\begin{enumerate}
  \item Hitunglah \omterm{def:g10:universal-gravitation:weight}{beratnya} di Bumi.
  \item Ia pergi ke Bulan. Berapa massanya dan berapa \omterm{def:g10:universal-gravitation:weight}{beratnya} di sana?
\end{enumerate}
\end{exercise}

\begin{exercise}[$\star$]\label{exo:g10:universal-gravitation:3}
\begin{enumerate}
  \item Dari kartu data, hitunglah kembali \omterm{def:g10:universal-gravitation:fieldstrength}{kuat medan gravitasi} di
        permukaan Mars.
  \item Sebuah penjelajah bermassa \qty{900}{kg} dikirim ke sana.
        Hitunglah \omterm{def:g10:universal-gravitation:weight}{beratnya} di Mars dan \omterm{def:g10:universal-gravitation:weight}{beratnya} di Bumi.
\end{enumerate}
\end{exercise}

\begin{exercise}[$\star$]\label{exo:g10:universal-gravitation:4}
Dengan memakai kartu data, hitunglah gaya yang dikerjakan Bumi pada
Bulan. Berapa gaya yang dikerjakan Bulan pada Bumi?
\end{exercise}

\begin{exercise}[$\star$]\label{exo:g10:universal-gravitation:5}
Dua benda saling menarik dengan gaya $F_0$. Menjadi berapa gayanya
apabila:
\begin{enumerate}
  \item jarak antara pusatnya digandakan;
  \item jaraknya dijadikan setengah;
  \item kedua massanya digandakan (jaraknya tetap);
  \item salah satu massanya dijadikan tiga kali lipat \emph{dan}
        jaraknya juga dijadikan tiga kali lipat?
\end{enumerate}
\end{exercise}

\begin{exercise}[$\star\star$]\label{exo:g10:universal-gravitation:6}
Dua kapal tanker raksasa bermuatan, masing-masing \qty{5.0e8}{kg},
ditambatkan dengan pusatnya terpisah \qty{100}{m}.
\begin{enumerate}
  \item Hitunglah tarikan \omterm{def:g10:universal-gravitation:force}{gravitasi} timbal balik keduanya.
  \item Massa berapa yang beratnya sama dengan gaya itu di Bumi?
  \item Gayanya besar secara mengejutkan --- namun kedua tanker itu
        sama sekali tidak menunjukkan tanda hendak saling mendekat.
        Hitunglah percepatan $a = F/m$ yang diberikannya kepada satu
        tanker, lalu berikan ulasanmu.
\end{enumerate}
\end{exercise}

\begin{exercise}[$\star\star$]\label{exo:g10:universal-gravitation:7}
Sebuah toko di Bulan menjual sekantong beras \qty{3.0}{kg}.
\begin{enumerate}
  \item Apa yang ditunjukkan neraca lengan dan neraca pegas untuk
        kantong ini di Bumi? Dan di Bulan?
  \item Seorang pembeli membayar untuk ``\qty{3.0}{kg}''. Alat manakah
        yang menjamin urusan yang adil di dunia mana pun, dan mengapa?
\end{enumerate}
\end{exercise}

\begin{exercise}[$\star\star$]\label{exo:g10:universal-gravitation:8}
Stasiun antariksa internasional terbang pada ketinggian \qty{400}{km}.
\begin{enumerate}
  \item Hitunglah \omterm{def:g10:universal-gravitation:fieldstrength}{kuat medan gravitasi} Bumi pada ketinggian itu.
        (Hati-hati: berapa jaraknya ke \emph{pusat} Bumi?)
  \item Nyatakan sebagai persentase nilainya di permukaan.
  \item Antariksawan di dalamnya melayang. Damaikan kenyataan itu
        dengan jawabanmu.
\end{enumerate}
\end{exercise}

\begin{exercise}[$\star\star$]\label{exo:g10:universal-gravitation:9}
Di permukaan Venus, $g = \qty{8.87}{N/kg}$, dan jari-jari planet itu
\qty{6.05e6}{m}. Simpulkan massa Venus dan bandingkan dengan massa
Bumi.
\end{exercise}

\begin{exercise}[$\star\star$]\label{exo:g10:universal-gravitation:10}
\begin{enumerate}
  \item Hitunglah gaya yang dikerjakan Matahari pada \qty{1.0}{kg} air
        laut, lalu gaya yang dikerjakan Bulan pada \omterm{def:g10:orders-of-magnitude:unit}{kilogram} yang sama.
  \item Matahari menarik lebih kuat --- dengan faktor berapa?
  \item Namun Bulanlah yang menguasai pasang surut. Pasang surut
        digerakkan bukan oleh tarikannya sendiri melainkan oleh
        \emph{selisih} antara tarikan pada sisi Bumi yang dekat dan
        yang jauh. Jelaskan, tanpa menghitung, mengapa Bulan yang dekat
        dapat mengalahkan Matahari yang jauh menurut ukuran itu.
\end{enumerate}
\end{exercise}

\begin{exercise}[$\star\star$]\label{exo:g10:universal-gravitation:11}
\begin{enumerate}
  \item Sebuah planet bermassa dua kali massa Bumi dan berjari-jari dua
        kali jari-jari Bumi. Nyatakan kuat medan di permukaannya dalam
        $g$ milik Bumi, lalu dalam \unit{N/kg}.
  \item Jari-jari berapa yang diperlukan sebuah planet bermassa sama
        dengan Bumi agar kuat medan di permukaannya menjadi $2g$?
\end{enumerate}
\end{exercise}

\begin{exercise}[$\star\star\star$]\label{exo:g10:universal-gravitation:12}
Di suatu tempat pada garis Bumi--Bulan, kedua tarikan pada sebuah wahana
antariksa saling meniadakan. Misalkan $x$ adalah jarak dari pusat Bumi
ke titik itu dan $D$ adalah jarak Bumi--Bulan.
\begin{enumerate}
  \item Tunjukkan bahwa di titik itu
        $\left(\dfrac{x}{D - x}\right)^{2}
        = \dfrac{M_{\text{Bumi}}}{M_{\text{Bulan}}}$.
  \item Simpulkan $x$. Berapa bagian dari perjalanan ke Bulan itu?
\end{enumerate}
\end{exercise}

\begin{exercise}[$\star\star\star$]\label{exo:g10:universal-gravitation:13}
Meriam Newton, dengan angka. Sebuah peluru menyerempet Bumi secara
mendatar dengan laju $v$.
\begin{enumerate}
  \item Setelah menempuh $x = \qty{8.0}{km}$ secara mendatar, Bumi yang
        bulat sudah ``turun menjauh'' sejauh $h$ yang memenuhi
        $(R + h)^2 = R^2 + x^2$. Dengan memakai
        $(R + h)^2 \approx R^2 + 2Rh$ untuk $h$ yang kecil, tunjukkan
        bahwa $h \approx x^2 / (2R)$ lalu hitunglah $h$.
  \item Benda yang jatuh bebas turun \qty{4.9}{m} pada sekon
        pertamanya. Bandingkan dengan $h$ lalu simpulkan laju yang
        membuat peluru itu tidak pernah mendarat.
  \item Hitunglah lama satu \omterm{def:g10:universal-gravitation:satellite}{orbit} penuh pada laju itu, sambil
        menyerempet permukaan.
\end{enumerate}
\end{exercise}

\begin{exercise}[$\star\star\star$]\label{exo:g10:universal-gravitation:14}
Sebuah bintang neutron memampatkan $M = \qty{2.8e30}{kg}$ (1.4 kali
massa Matahari) ke dalam jari-jari \qty{12}{km}.
\begin{enumerate}
  \item Hitunglah \omterm{def:g10:universal-gravitation:fieldstrength}{kuat medan gravitasi} di permukaannya.
  \item Berapa \omterm{def:g10:universal-gravitation:weight}{berat} manusia \qty{70}{kg} di sana? Bandingkan dengan
        \omterm{def:g10:universal-gravitation:weight}{beratnya} di Bumi.
  \item Berapa \omterm{def:g10:universal-gravitation:weight}{berat} sebutir pasir \qty{1.0}{mg} di sana? Massa berapa
        yang memiliki \omterm{def:g10:universal-gravitation:weight}{berat} sebesar itu di Bumi?
\end{enumerate}
\end{exercise}

\begin{exercise}[$\star\star\star$]\label{exo:g10:universal-gravitation:15}
Uji Bulan yang dilakukan Newton (1666). Jari-jari \omterm{def:g10:universal-gravitation:satellite}{orbit} Bulan kira-kira
$60$ jari-jari Bumi.
\begin{enumerate}
  \item Dengan faktor berapa tarikan Bumi per \omterm{def:g10:orders-of-magnitude:unit}{kilogram} lebih lemah di
        Bulan daripada di permukaan Bumi, jika \omterm{def:g10:universal-gravitation:force}{gravitasi} mengikuti
        kebalikan kuadrat?
  \item Di dekat permukaan, benda yang jatuh turun \qty{4.9}{m} pada
        sekon pertamanya. Seberapa jauh Bulan seharusnya jatuh ke arah
        Bumi tiap sekon?
  \item Periksalah dengan \omterm{def:g10:universal-gravitation:satellite}{orbitnya}: Bulan menempuh lingkarannya yang
        berjari-jari $d = \qty{3.84e8}{m}$ dalam $T = 27.3$ hari.
        Hitunglah lajunya $v$, jarak $x$ yang ditempuhnya dalam satu
        sekon, lalu jatuhnya $h \approx x^2/(2d)$ (seperti pada
        \cref{exo:g10:universal-gravitation:13}).
  \item Bandingkan jawaban pertanyaan 2 dan 3, lalu nyatakan apa yang
        disimpulkan Newton.
\end{enumerate}
\end{exercise}

\section{Soal: Menimbang Bumi}

\begin{problem}[{Soal akhir pekan --- dari dua bola timbal di sebuah
gudang sampai ke massa Bumi lalu massa Matahari: apa gunanya satu
tetapan kecil yang diukur sekali saja}]
\label{pb:g10:universal-gravitation:1}
Pada tahun 1798 Henry Cavendish menggantungkan sebatang ringan yang
membawa dua bola timbal kecil pada seutas kawat tipis, mendekatkan dua
bola timbal besar, lalu mengukur puntiran kawat yang nyaris tidak ada.
Surat kabar mengatakan ia telah \emph{menimbang Bumi} --- dan mereka
benar: begitu $G$ diketahui, $g$ dan jari-jari Bumi sendiri menyerahkan
massanya, dan \omterm{def:g10:universal-gravitation:satellite}{orbit} tahunan Bumi menyerahkan massa Matahari. Soal ini
menelusuri kembali seluruh rantai itu. Pakailah kartu data di sepanjang
soal.

\medskip
\noindent\textbf{Bagian I --- Gaya yang paling lemah.}
\begin{enumerate}
  \item Besi Menara Eiffel bermassa sekitar \qty{7.3e6}{kg}. Seorang
        pelancong bermassa \qty{55}{kg} berdiri \qty{100}{m} dari pusat
        massanya. Hitunglah tarikan menara itu pada si pelancong.
  \item Bandingkan dengan \omterm{def:g10:universal-gravitation:weight}{berat} sebutir pasir \qty{1.0}{mg}.
  \item Latihan kesebandingan: apa yang terjadi pada \omterm{def:g10:universal-gravitation:force}{gaya gravitasi}
        apabila (a) jaraknya dikalikan $10$; (b) kedua massanya
        dikalikan $1000$; (c) kedua massanya \emph{dan} jaraknya
        dikalikan $1000$?
  \item Hitunglah tarikan Bumi pada benda \qty{1.00}{kg} di
        permukaannya, dari $M$, $R$ dan $G$. Nama dan lambang
        sehari-hari apa yang dibawa bilangan ini?
  \item \omterm{def:g10:universal-gravitation:force}{Gravitasi} jauh merupakan interaksi terlemah yang pernah kamu
        temui --- namun gravitasilah yang menjalankan alam semesta.
        Jelaskan dalam satu atau dua kalimat mengapa tarikan Bumi
        padamu mengalahkan tarikan tetanggamu.
\end{enumerate}

\medskip
\noindent\textbf{Bagian II --- Neraca Cavendish.}
Di dalam neraca puntir itu, sebuah bola besar bermassa
$m_1 = \qty{158}{kg}$ menarik bola kecil bermassa
$m_2 = \qty{0.73}{kg}$; pusat keduanya terpisah $d = \qty{0.23}{m}$.
\begin{enumerate}[resume]
  \item Dengan nilai $G$ yang modern, hitunglah gaya antara kedua bola
        itu.
  \item Hitunglah \omterm{def:g10:universal-gravitation:weight}{berat} bola kecilnya dan nisbah kedua gayanya. Mengapa
        Cavendish memerlukan kawat puntir dan bukan neraca biasa?
  \item Puntiran yang diukur Cavendish bersesuaian (dalam \omterm{def:g10:orders-of-magnitude:unit}{satuan}
        modern) dengan $F = \qty{1.47e-7}{N}$. Simpulkan nilai $G$
        miliknya dan bandingkan dengan nilai masa kini.
  \item Dengan $G = \qty{6.67e-11}{N.m^2/kg^2}$, $g = \qty{9.81}{N/kg}$
        dan $R = \qty{6.37e6}{m}$, simpulkan massa Bumi. (Inilah
        langkah yang dirayakan surat kabar itu.)
  \item Simpulkan massa jenis rata-rata Bumi $\rho = M / V$, dengan
        $V = \frac43 \pi R^3$. Batuan permukaan bermassa jenis sekitar
        \qty{2.7e3}{kg/m^3}: apa yang diungkapkan perbandingan itu
        tentang bagian dalam Bumi yang jauh di bawah?
\end{enumerate}

\medskip
\noindent\textbf{Bagian III --- Dunia yang lain.}
\begin{enumerate}[resume]
  \item Dari kartu data, hitunglah kembali kuat medan di permukaan
        Bulan.
  \item Di Bumi, tolakan kaki yang baik mengangkat pusat massamu
        sejauh \qty{0.50}{m}. Untuk tolakan yang sama, tinggi yang
        dicapai berbanding terbalik dengan $g$ (sebagaimana akan
        dibenarkan bab energi, \cref{ch:g10:energy-conservation}).
        Setinggi apa lompatan yang sama membawamu di Bulan?
  \item Merkurius: $M = \qty{3.30e23}{kg}$, $R = \qty{2.44e6}{m}$.
        Hitunglah kuat medan di permukaannya dan bandingkan dengan
        Mars. Merkurius jauh lebih kecil daripada Mars --- bagaimana
        kedua nilainya bisa sedemikian dekat?
  \item Pada ketinggian berapa di atas permukaan Bumi kuat medannya
        turun menjadi setengah nilainya di permukaan?
  \item Pada jarak berapa dari pusat Bumi tarikan Bumi per \omterm{def:g10:orders-of-magnitude:unit}{kilogram}
        turun sampai nilai permukaan Bulan, yaitu \qty{1.62}{N/kg}?
        Nyatakan dalam jari-jari Bumi.
\end{enumerate}

\medskip
\noindent\textbf{Bagian IV --- Menimbang Matahari.}
Bumi menempuh \omterm{def:g10:universal-gravitation:satellite}{orbit} yang hampir berbentuk lingkaran berjari-jari
$d = \qty{1.496e11}{m}$ dalam satu tahun, $T = \qty{3.156e7}{s}$.
Terimalah sebagai diketahui (hal ini dibuktikan pada
\cref{ch:g12:satellites-kepler}) bahwa benda pada lintasan lingkaran
berjari-jari $d$ dengan laju $v$ memerlukan tarikan sebesar $v^2/d$
newton per \omterm{def:g10:orders-of-magnitude:unit}{kilogram} menuju pusatnya.
\begin{enumerate}[resume]
  \item Hitunglah laju \omterm{def:g10:universal-gravitation:satellite}{orbit} Bumi $v$.
  \item Simpulkan tarikan per \omterm{def:g10:orders-of-magnitude:unit}{kilogram} yang dikerjakan Matahari pada
        Bumi.
  \item Tarikan itu juga sama dengan $G M_{\text{Matahari}} / d^2$.
        Simpulkan massa Matahari.
  \item Berapa banyak Bumi itu? Bandingkan $M_{\text{Matahari}}$
        dengan $M_{\text{Bumi}}$.
  \item Penutup. Bola-bola Cavendish, Bumi, Matahari: satu tetapan $G$
        melayani ketiganya. Nyatakan dalam satu kalimat apa yang
        dibelikan kata \emph{semesta} pada ``\omterm{def:g10:universal-gravitation:force}{gravitasi} semesta''
        kepada kita di dalam soal ini.
\end{enumerate}
\end{problem}
